\clearpage
\begin{examplebox}
\textbf{\textcolor{CExample}{例1：已知一段，求另一段}}

\textbf{题目：}$y^2=4x$，焦点弦满足 $AF=3$，求 $BF$。
\begin{center}\begin{tikzpicture}[font=\small,>=stealth]
\draw[rounded corners=3pt,gray!45] (0,0) rectangle (14.2,11.5);
\node at (7.1,10.95) {已知一段，另一段怎样变化？先看长度配对的曲线};

\begin{scope}[shift={(1.1,2.2)},scale=.95]
\draw[->,gray] (0,0)--(7.8,0) node[below=7pt] {$r=AF$};
\draw[->,gray] (0,0)--(0,7.8) node[above right] {$s=BF$};
\draw[violet,densely dashed] (1,0)--(1,7.5);
\draw[violet,densely dashed] (0,1)--(7.5,1);
\draw[teal!75!black,very thick,samples=150,domain=1.155:7.5,variable=\t] plot ({\t},{\t/(\t-1)});
\node[below=5pt] at (0,0) {$0$};
\node[below=5pt,violet] at (1,0) {$1$};
\node[left=5pt,violet] at (0,1) {$1$};
\fill[blue!75!black] (2,2) circle(2.5pt);
\draw[blue!75!black,->] (3.3,4.4)--(2.05,2.1);
\node[blue!75!black,above,align=center] at (3.3,4.4) {通径\\$r=s=2$};
\fill[orange!85!black] (3,1.5) circle(2.5pt);
\draw[orange!85!black,->] (5.45,3.2)--(3.1,1.55);
\node[orange!85!black,above,align=center] at (5.45,3.2) {例1\\$r=3,\ s=3/2$};
\end{scope}
\node[anchor=west,align=left] at (9.3,6.6)
 {$p=2$，$\displaystyle s=\frac r{r-1}$\\[14pt]
 $r>1,\quad s>1$\\[12pt]
 一段变长，另一段变短\\[12pt]
 $r=1$ 没有有限配对\\[10pt]
 虚线只可逼近，不可取到};
\node at (7.1,.8) {图中每个点是一对长度 $(r,s)$，不是抛物线上的端点坐标。};
\end{tikzpicture}
\end{center}
由 $2p=4$ 得 $p=2$。令 $BF=s$，则
\[
 \frac13+\frac1s=1\quad\Longrightarrow\quad
 \frac1s=\frac23\quad\Longrightarrow\quad\boxed{s=\frac32}.
\]
图中橙点就是配对 $(3,3/2)$。回验 $1/3+2/3=1$，且两段均大于 $p/2=1$。
\end{examplebox}

\clearpage
\begin{examplebox}
\textbf{\textcolor{CExample}{例2：给倾角，分别求倒数和与弦长}}

\textbf{题目：}$y^2=8x$ 的焦点弦倾斜角为 $60^\circ$，求倒数和及 $AB$。
由 $p=4$，倒数和为 $\boxed{2/p=1/2}$，与角度无关。
按上、下端标号，
\[
 AF=\frac4{1-1/2}=8,\quad BF=\frac4{1+1/2}=\frac83,
 \quad\boxed{AB=\frac{32}3}.
\]
回验 $1/8+3/8=1/2$，且 $AB=2p/\sin^260^\circ=32/3$。

\par\medskip
\textbf{\textcolor{CExample}{例3：给比例，只设一个未知数}}

\textbf{题目：}$y^2=4x$ 的焦点弦满足 $AF:BF=3:1$，求倒数和及 $AB$。
$p=2$，故倒数和为 $1$。设 $AF=3t$、$BF=t$，$t>0$，则
\[
 \frac1{3t}+\frac1t=1\quad\Longrightarrow\quad t=\frac43.
\]
所以 $\boxed{AF=4,\ BF=4/3,\ AB=16/3}$。
这正是前面相似图的两段；回验 $1/4+3/4=1$。

\par\medskip
\textbf{\textcolor{CExample}{迁移训练：比例换成 $2:3$}}

\textbf{题目：}$y^2=12x$ 的焦点弦满足 $AF:BF=2:3$，求两段及弦长。
$p=6$，设 $AF=2t$、$BF=3t$，则
\[
 \frac1{2t}+\frac1{3t}=\frac13
 \quad\Longrightarrow\quad t=\frac52.
\]
故 $\boxed{AF=5,\ BF=15/2,\ AB=25/2}$。
回验 $1/5+2/15=1/3=2/p$。
\end{examplebox}

\clearpage
\begin{examplebox}
\textbf{\textcolor{CExample}{例4：由倒数和推出通径最短}}

\textbf{题目：}对 $y^2=6x$，证明通径最短，并求其长度。
由 $p=3$，记 $r=AF,s=BF$，则 $1/r+1/s=2/3$。
\[
 (r+s)\left(\frac1r+\frac1s\right)
 =2+\frac rs+\frac sr\ge4.
\]
故 $\boxed{r+s\ge4/(2/3)=6}$。
等号要求 $r=s$，再由倒数和得 $r=s=3$；这正是通径。
\textbf{回验：}竖直线 $x=3/2$ 的端点为 $(3/2,\pm3)$，弦长为 $6$。

\par\medskip
\textbf{\textcolor{CExample}{例5：检查“已知长度”是否可能}}

\textbf{题目：}$y^2=4x$ 是否存在有限焦点弦使 $AF=1$？
此时 $p=2$。若 $AF=1$，则
\[
 1+\frac1{BF}=1\quad\Longrightarrow\quad\frac1{BF}=0,
\]
不可能由有限的正长度 $BF$ 满足，所以\textbf{不存在}。
一般地，$1/s=2/p-1/r>0$，因此
\[
 \boxed{r>p/2},\qquad \boxed{s>p/2}.
\]
回看例1配对曲线的两条虚线：长度可以无限接近下界，不能取到。
注意 $p/2$ 是\textbf{严格下界}，不是可取得的最小值。
\end{examplebox}
